\PassOptionsToPackage{quiet}{xeCJK}
\documentclass[answer]{THUExam}
\usepackage{HTNotes-math}
\usepackage{pifont}
\usepackage{ulem}

\usepackage{tikz}
\usetikzlibrary{positioning, calc, shapes.geometric, shapes, shapes.multipart, arrows.meta, arrows, decorations.markings, external, trees}

% =================================================
%       PDF信息
% =================================================
\title{高等数学A (二)}
\author{}
\Subject{作业}
\Keywords{高等数学, 多元函数微分学, 几何应用}

% =================================================
%       试卷头信息
% =================================================
\Year{2025--2026}
\Semester{春季}
\Course{高等数学A}
\Time{}
\Type{A}

\begin{document}
\vspace{1cm}
\makehead

\vspace{1cm}

% =================================================
%       选择题
% =================================================
\makepart{选择题}{每小题~4~分，共~16~分}

\begin{problem}
函数 $f(x,y)$ 在点 $(x_0,y_0)$ 处沿方向 $\vec{l}$ 的方向导数取最大值时，$\vec{l}$ 的方向为\pickout{B}
\options
{任意方向}
{梯度的方向}
{梯度的反方向}
{等值线的切线方向}
\end{problem}
\begin{note}
方向导数 $f_{\vec{l}}=\nabla f\cdot\vec{l}^{\kern1pt0}$ 在内积取最大值时等于 $|\nabla f|$，此时单位向量 $\vec{l}^{\kern1pt0}$ 与梯度 $\nabla f$ 同向。
\end{note}
\bigskip

\begin{problem}
曲面 $F(x,y,z)=0$ 在点 $(x_0,y_0,z_0)$ 处的一个法向量为\pickout{A}
\options
{$(F_x,F_y,F_z)\big|_{P_0}$}
{$(F_x,F_y,-1)\big|_{P_0}$}
{$(z_x,z_y,-1)\big|_{P_0}$}
{$(1,1,1)$}
\end{problem}
\begin{note}
曲面 $F(x,y,z)=0$ 在一点的梯度 $\nabla F=(F_x,F_y,F_z)$ 即为曲面在该点的一个法向量。
\end{note}
\bigskip

\begin{problem}
空间曲线 $\varGamma:x=x(t),\,y=y(t),\,z=z(t)$ 在点 $t=t_0$ 处的切向量可取为\pickout{B}
\options
{$(x(t_0),y(t_0),z(t_0))$}
{$(x'(t_0),y'(t_0),z'(t_0))$}
{$(x''(t_0),y''(t_0),z''(t_0))$}
{$(x'(t_0),0,0)$}
\end{problem}
\begin{note}
空间曲线以参数方程给出时，切向量 $\vec{T}$ 的分量为各坐标函数对参数的导数，即 $\vec{T}=(x'(t),y'(t),z'(t))$。
\end{note}
\bigskip

\begin{problem}
设 $f(x,y)$ 在点 $P_0$ 可微，$\vec{l}^{\kern1pt0}$ 为给定的单位方向向量，则方向导数 $\dfrac{\partial f}{\partial l}\Big|_{P_0}$ 的计算公式为\pickout{C}
\options
{$\nabla f(P_0)\cdot\vec{l}$，其中 $\vec{l}$ 为任意方向向量}
{$|\nabla f(P_0)|\cos\theta$}
{$\nabla f(P_0)\cdot\vec{l}^{\kern1pt0}$}
{$\nabla f(P_0)\times\vec{l}^{\kern1pt0}$}
\end{problem}
\begin{note}
方向导数定义为 $f_{\vec{l}}=\nabla f\cdot\vec{l}^{\kern1pt0}$，其中 $\vec{l}^{\kern1pt0}$ 必须是单位向量。若方向向量 $\vec{l}$ 非单位，需先将其单位化。
\end{note}
\bigskip

% =================================================
%       简答题
% =================================================
\makepart{简答题}{共~3~题，满分~34~分}

\begin{problem}
求空间曲线
\[
x=t,\quad y=t^2,\quad z=t^3
\]
在 $t=1$ 处的切线方程和法平面方程。
\end{problem}

\begin{solution}
切点坐标：$(1,1,1)$。

求导得切向量：
\[
\vec{T}=(x'(t),y'(t),z'(t))=(1,\,2t,\,3t^2)\;\xrightarrow{t=1}\;(1,2,3).
\]

切线方程（点向式）：
\[
\frac{x-1}{1}=\frac{y-1}{2}=\frac{z-1}{3}.
\]

法平面方程（切向量即法平面的法向量）：
\[
1\cdot(x-1)+2\cdot(y-1)+3\cdot(z-1)=0
\;\Longrightarrow\; x+2y+3z=6.
\]
\end{solution}
\bigskip

\begin{problem}
求曲面 $z=x^2+y^2$ 在点 $(1,1,2)$ 处的切平面方程和法线方程。
\end{problem}

\begin{solution}
令 $F(x,y,z)=x^2+y^2-z=0$，则
\[
F_x=2x,\quad F_y=2y,\quad F_z=-1.
\]

在点 $(1,1,2)$ 处，法向量
\[
\vec{n}=(F_x,F_y,F_z)\big|_{(1,1,2)}=(2,2,-1).
\]

切平面方程：
\[
2(x-1)+2(y-1)-(z-2)=0 \;\Longrightarrow\; 2x+2y-z=2.
\]

法线方程（点向式）：
\[
\frac{x-1}{2}=\frac{y-1}{2}=\frac{z-2}{-1}.
\]
\end{solution}
\bigskip

\begin{problem}
求函数 $f(x,y,z)=xy+yz+zx$ 在点 $P_0(1,1,1)$ 处的梯度，并求函数在该点沿方向 $\vec{l}=(1,1,1)$ 的方向导数。
\end{problem}

\begin{solution}
先求偏导数：
\[
f_x=y+z,\quad f_y=x+z,\quad f_z=x+y.
\]

在 $P_0(1,1,1)$ 处：
\[
\nabla f(P_0)=(f_x,f_y,f_z)\big|_{P_0}=(2,2,2).
\]

$\vec{l}=(1,1,1)$ 的单位向量：
\[
\vec{l}^{\kern1pt0}=\frac{(1,1,1)}{\sqrt{1^2+1^2+1^2}}
=\Bigl(\frac{1}{\sqrt{3}},\;\frac{1}{\sqrt{3}},\;\frac{1}{\sqrt{3}}\Bigr).
\]

方向导数：
\[
\frac{\partial f}{\partial l}\Big|_{P_0}
=\nabla f(P_0)\cdot\vec{l}^{\kern1pt0}
=2\cdot\frac{1}{\sqrt{3}}+2\cdot\frac{1}{\sqrt{3}}+2\cdot\frac{1}{\sqrt{3}}
=\frac{6}{\sqrt{3}}=2\sqrt{3}.
\]
\end{solution}
\bigskip

\end{document}
